




\documentclass[a4paper,12pt]{article}
\usepackage[utf8]{inputenc}
\usepackage[T2A]{fontenc}
\usepackage[russian]{babel}
\usepackage{amsmath}
\usepackage{graphicx}
\usepackage{geometry}
\usepackage{float}
\usepackage{booktabs}
\usepackage{caption}
\usepackage{subcaption}
\usepackage{multirow}
\usepackage{wrapfig}
\geometry{top=2cm, bottom=2cm, left=2.5cm, right=1.5cm}

\title{ОТЧЕТ \\ ПО ЛАБОРАТОРНОЙ РАБОТЕ \\ <<Атомно-силовой микроскоп>>}
\author{
    Иванина Екатерина, Пенская Дарья, Питерская Екатерина \\
    \vspace{0.5cm} \\
}
\date{25 ноября 2025 г.}

\begin{document}

% Титульный лист
\begin{titlepage}
    \centering
    \vspace*{1cm}

    \textbf{Министерство науки и высшего образования Российской Федерации} \\
    \vspace{0.5cm}
    \textbf{ФЕДЕРАЛЬНОЕ ГОСУДАРСТВЕННОЕ АВТОНОМНОЕ} \\
    \textbf{ОБРАЗОВАТЕЛЬНОЕ УЧРЕЖДЕНИЕ ВЫСШЕГО ОБРАЗОВАНИЯ} \\
    \vspace{0.5cm}
    \textbf{«МОСКОВСКИЙ ФИЗИКО-ТЕХНИЧЕСКИЙ ИНСТИТУТ} \\
    \textbf{(НАЦИОНАЛЬНЫЙ ИССЛЕДОВАТЕЛЬСКИЙ УНИВЕРСИТЕТ)»} \\
    \textbf{(МФТИ, Физтех)} \\
    \vspace{1cm}
    \textbf{КАФЕДРА ВАКУУМНОЙ ЭЛЕКТРОНИКИ} \\

    \vspace{2cm}

    \textbf{ОТЧЕТ} \\
    \textbf{ПО ЛАБОРАТОРНОЙ РАБОТЕ} \\
    \vspace{1cm}
    \textbf{\Large <<Атомно-силовой микроскоп>>} \\

    \vspace{2cm}

    \begin{flushright}
        Работу выполнили: \\
        Иванина Екатерина \\
        Пенская Дарья \\
        Питерская Екатерина \\
       
        \vspace{1cm}
        Работу принял, оценка: \underline{\hspace{3cm}} \\
        \vspace{0.5cm}
        (подпись, дата, оценка)
    \end{flushright}

    \vfill
    Долгопрудный \\ 2025
\end{titlepage}

% Аннотация
\section*{Аннотация}
\textbf{Цель работы:} познакомиться с физическими принципами работы атомно-силового микроскопа (АСМ), исследовать топографию игольчатой подложки (тестовой решетки TGT1) и CD-диска, оценить радиус кривизны зонда.

\section{Теоретическая часть}

\subsection{Физические принципы работы АСМ}

Атомно-силовая микроскопия (АСМ) основана на регистрации силового взаимодействия между зондом и поверхностью исследуемого образца. Характер взаимодействия определяется свойствами зонда, образца и окружающей среды.

\subsubsection{Типы силовых взаимодействий}

Различают короткодействующие и дальнодействующие силы. 

\textbf{Короткодействующие силы} возникают при перекрытии электронных оболочек атомов иглы и поверхности на расстояниях до нескольких ангстрем. Эти силы обеспечивают атомарное разрешение и препятствуют проникновению иглы в образец.

\textbf{Дальнодействующие силы} включают:
\begin{itemize}
    \item Ван-дер-ваальсовы силы ($F_{ВдВ} \approx 10^{-9}\div10^{-8}$ Н)
    \item Капиллярные силы ($F_{кап} \approx 10^{-8}$ Н)
    \item Электростатические и магнитные взаимодействия
\end{itemize}
Они проявляются на расстояниях до тысяч ангстрем.

\subsubsection{Режимы работы АСМ}

\begin{figure}[H]
    \centering
    \includegraphics[width=0.5\linewidth]{режимы.png}
    \caption{Основные режимы работы АСМ}
    \label{fig:acm_modes}
\end{figure}

В АСМ используются три основных режима работы:
\begin{itemize}
    \item \textbf{Контактный} -- постоянное касание поверхности
    \item \textbf{Бесконтактный} -- зонд на расстоянии нескольких нанометров
    \item \textbf{Полуконтактный} -- периодическое касание поверхности
\end{itemize}

\subsubsection{Преимущества АСМ}

Ключевые преимущества АСМ:
\begin{itemize}
    \item Возможность исследования непроводящих материалов
    \item Нанометровое и атомарное разрешение
    \item Работа в атмосферных условиях
    \item Изучение различных свойств поверхности: магнитных и электрических полей, температурных распределений, ёмкостных и адгезионных свойств
    \item Исследование процессов на границе жидкость-твёрдое тело
\end{itemize}

Для достижения атомарного разрешения необходимы чистая поверхность, защита от вибраций и проведение измерений в вакууме.

\subsection{Устройство атомно-силового микроскопа}

\subsubsection{Общая схема АСМ}

Атомно-силовой микроскоп состоит из четырех основных компонентов:
\begin{itemize}
    \item Зонд
    \item Система регистрации отклонения зонда
    \item Пьезосканер
    \item Система обратной связи
\end{itemize}

\begin{figure}[H]
    \centering
    \includegraphics[width=0.7\linewidth]{схема-асм.png}
    \caption{Схема взаимодействия основных элементов АСМ}
    \label{fig:acm_scheme}
\end{figure}

Сканирование поверхности осуществляется пьезосканером, который перемещает образец относительно зонда по координатам X и Y. Взаимодействие зонда с поверхностью вызывает отклонение кантилевера, которое регистрируется измерительной системой.

\subsubsection{Режимы сканирования}

В атомно-силовой микроскопии используются три основных режима сканирования:

\begin{itemize}
    \item \textbf{Режим постоянной высоты} -- образец движется при постоянной координате $Z$, измеряется сила взаимодействия $F_{ij}(Z = const)$
    \item \textbf{Режим постоянного взаимодействия} -- система обратной связи поддерживает постоянную силу $F(X, Y, Z)$ за счет изменения координаты $Z$
    \item \textbf{Режим спектроскопии} -- при $X, Y = const$ измеряется зависимость сигнала $S_i$ от параметра взаимодействия $p_i$
\end{itemize}

\subsubsection{Конструкция зонда}

\begin{figure}[H]
    \centering
    \includegraphics[width=0.8\linewidth]{зонд.png}
    \caption{Конструкция зонда АСМ}
    \label{fig:probe_structure}
\end{figure}

Зонд АСМ состоит из кантилевера - балки с закрепленным и свободным концами, на свободном конце которого расположена острая игла.

Основные параметры иглы:
\begin{itemize}
    \item Радиус кривизны менее 10 нм
    \item Длина 3-15 мкм
    \item Угол при вершине: около 22°
\end{itemize}

Кантилеверы изготавливаются из кремния или нитрида кремния и могут иметь треугольную (V-образную) или прямоугольную форму длиной 80-350 мкм. Для обеспечения чувствительности к сверхмалым силам жёсткость кантилевера должна быть минимальной (0.01-100 Н/м), а резонансная частота - максимально высокой (15-500 кГц).

\subsubsection{Система регистрации отклонений}

\begin{figure}[H]
    \centering
    \begin{subfigure}{0.45\textwidth}
        \centering
        \includegraphics[width=\linewidth]{фигня1.png}
        \caption{Оптическая схема}
    \end{subfigure}
    \hfill
    \begin{subfigure}{0.45\textwidth}
        \centering
        \includegraphics[width=\linewidth]{фигня2.png}
        \caption{Фотодиодная система}
    \end{subfigure}
    \caption{Система регистрации отклонений кантилевера}
    \label{fig:detection_system}
\end{figure}

Наиболее распространённым методом регистрации отклонений кантилевера является оптический метод. Лазерный луч с длиной волны 670 нм и мощностью 0.9 мВт фокусируется на кантилевер, а отражённый луч попадает на четырехсекционный фотодиод.

Система формирует три сигнала:
\begin{itemize}
    \item \textbf{DFL} = (A + B) - (C + D) -- вертикальное отклонение
    \item \textbf{LF} = (A + C) - (B + D) -- боковое отклонение
    \item \textbf{LASER} = A + B + C + D -- интенсивность отражения
\end{itemize}

Альтернативным методом является использование тензорезистивных кантилеверов, основанных на изменении электросопротивления при механической деформации.

\begin{figure}[H]
    \centering
    \includegraphics[width=0.8\linewidth]{фигня3.png}
    \caption{Тензорезистивный кантилевер}
    \label{fig:strain_cantilever}
\end{figure}

\subsubsection{Пьезосканер и его характеристики}

Пьезосканер обеспечивает два независимых движения образца относительно кантилевера: сканирование вдоль поверхности образца и перемещение в направлении, перпендикулярном поверхности. Изготовленный из цирконата-титаната свинца, сканер имеет пьезоэлектрический коэффициент от 0.1 до 300 нм/В.

Конструктивно сканеры могут выполняться в виде пьезотрубок (внутренняя для перемещения по X, Y, внешняя по Z) или на основе биморфных элементов. Поле сканирования варьируется от $3 \times 3$ мкм$^2$ до $100 \times 100$ мкм$^2$ при максимальной измеряемой высоте от 2 до 10 мкм.

\begin{figure}[H]
    \centering
    \includegraphics[width=0.5\linewidth]{фигня4.png}
    \caption{Пьезосканер на основе пьезотрубок}
    \label{fig:piezo_scanner}
\end{figure}

\subsubsection{Особенности пьезосканеров}

Пьезосканеры обладают рядом особенностей, влияющих на точность измерений:

\begin{itemize}
    \item \textbf{Нелинейность} -- пьезокерамика деформируется нелинейно с приложенным напряжением. Обусловлено уменьшением пьезомодуля на 10-20\% с ростом напряжения
    \item \textbf{Гистерезис} -- неоднозначная зависимость удлинения от направления изменения напряжения
    \item \textbf{Ползучесть} -- медленный дрейф после быстрого изменения напряжения
    \item \textbf{Температурный дрейф} -- изменение размеров при колебаниях температуры
\end{itemize}

\begin{figure}[H]
    \centering
    \includegraphics[width=0.4\linewidth]{фигня5.png}
    \caption{Гистерезис и ползучесть пьезокерамики}
    \label{fig:piezo_hysteresis}
\end{figure}

\subsubsection{Система обратной связи}

Система обратной связи на основе PID-регулятора поддерживает постоянное взаимодействие между зондом и образцом. Выходной сигнал регулятора FBO представляет сумму трех компонент:

\begin{itemize}
    \item $W_P = P(FBI(t) - SP)$ -- пропорциональная компонента
    \item $W_I = I \int_{t-r}^{t} (FBI(t') - SP) dt'$ -- интегральная компонента
    \item $W_D = D \frac{\partial (FBI(t) - SP)}{\partial t}$ -- дифференциальная компонента
\end{itemize}

Пропорциональная компонента обеспечивает отклик на резкие изменения, интегральная -- отрабатывает крупные детали поверхности, а дифференциальная -- демпфирует нежелательные осцилляции. Правильный подбор параметров P, I, D обеспечивает точное отслеживание рельефа поверхности.

\subsection{Методики измерения в атомно-силовой микроскопии}

Основные методики атомно-силовой микроскопии можно разделить на четыре группы: спектроскопия, однопроходное и многопроходное сканирование, а также зондовая наполитография.

\subsubsection{Спектроскопические методы}

В спектроскопических режимах измеряется зависимость различных сигналов от расстояния между зондом и образцом или от приложенного напряжения. Это даёт возможность исследовать локальные силы взаимодействия, механические и электрические свойства поверхности.

\subsubsection{Однопроходное сканирование}

Однопроходное сканирование применяется для получения стандартного топографического изображения построчно. При этом параметры обратной связи фиксированы, и карта рельефа формируется по изменению сигнала детектора.

\subsubsection{Многопроходные режимы}

Многопроходные режимы предполагают повторное сканирование одной и той же линии с разными настройками. Это расширяет набор измеряемых величин и позволяет получать карты электрических, магнитных и других свойств образца с учётом уже известного рельефа.

\subsubsection{Зондовая наполитография}

АСМ может использоваться для локального воздействия на поверхность -- механического или электрического -- что лежит в основе методов зондовой наполитографии.

\subsection{Режимы работы АСМ}

\subsubsection{Контактный режим работы АСМ}

Контактный режим является классическим способом работы АСМ и основан на поддержании постоянной силы взаимодействия между зондом и поверхностью. 

\textbf{Преимущества:}
\begin{itemize}
    \item Высокая скорость сканирования
    \item Возможность получения изображений с высоким пространственным разрешением, вплоть до атомного уровня
    \item Возможность измерения локальных сил трения и локальной жесткости
\end{itemize}

\textbf{Ограничения:}
\begin{itemize}
    \item Повышенное воздействие на поверхность
    \item Риск модификации мягких образцов
    \item Значительные силы трения и капиллярный эффект
\end{itemize}

\subsubsection{Полуконтактный режим работы АСМ}

Полуконтактный режим основан на возбуждении кантилевера на его резонансной частоте. Зонд периодически касается поверхности образца.

\textbf{Преимущества:}
\begin{itemize}
    \item Снижение механического воздействия на образец
    \item Повышенная чувствительность по сравнению со статическим измерением
    \item Уменьшение величины шумов за счет переноса спектра в область высоких частот
    \item Возможность одновременной регистрации амплитуды и фазы колебаний
\end{itemize}

\textbf{Ограничения:}
\begin{itemize}
    \item Зависимость сигналов от параметров резонанса
    \item Влияние условий внешней среды
\end{itemize}

\subsubsection{Бесконтактный режим работы АСМ}

В бесконтактном режиме кантилевер колеблется на резонансной частоте, не касаясь поверхности образца. Для детектирования используется модуляционная методика.

\textbf{Преимущества:}
\begin{itemize}
    \item Полное отсутствие контакта с поверхностью
    \item Отсутствие сил трения и износа зонда
    \item Подходит для особо мягких образцов
\end{itemize}

\textbf{Недостатки:}
\begin{itemize}
    \item Высокая чувствительность к вибрациям и внешней среде
    \item Зависимость сигнала от слабых дальнодействующих сил
    \item Обычно более низкое разрешение по сравнению с контактными режимами
\end{itemize}

\section{Результаты измерений и обработка данных}

В ходе работы были получены изображения структуры решётки TGT01 и CD-диска в контактном режиме работы. Для CD-диска также получены зависимости амплитуды колебаний зонда от расстояния зонд-образец.

Перед каждым изображением производилось нахождение резонансной частоты для корректной работы прибора. Пример амплитудно-частотной характеристики (АЧХ) зонда представлен на рис. \ref{fig:afm_characteristic}.

\begin{figure}[H]
    \centering
    \includegraphics[width=0.8\linewidth]{image.png}
    \caption{Амплитудно-частотная характеристика зонда до начала измерений}
    \label{fig:afm_characteristic}
\end{figure}

\subsection{Исследование CD-диска}

Установим на предметный столик АСМ образец из кусочка CD-диска. Проведём измерения амплитудно-частотной характеристики (АЧХ) зонда для определения резонансной частоты материала кантилевера.

Как отмечено в методике, резонансная частота собственных колебаний кантилевера должна быть максимально высокой, чтобы собственные колебания не возбуждались за счёт возвратно-поступательных перемещений иглы относительно образца. В данной работе для кантилевера требовалось установить амплитуду раскачки в диапазоне 8--12 мВ. 

\begin{figure}[H]
    \centering
    \includegraphics[width=0.8\linewidth]{image2.png}
    \caption{Амплитудно-частотная характеристика зонда для определения резонансной частоты}
    \label{fig:resonance_determination}
\end{figure}

Экспериментально получена резонансная частота колебаний зонда $f_0 = 7.71$ кГц.

\begin{figure}[H]
    \centering
    \includegraphics[width=0.8\linewidth]{image3.png}
    \caption{Полученное изображение поверхности CD-диска в 2D и 3D}
    \label{fig:cd_surface}
\end{figure}

\begin{figure}[H]
    \centering
    \begin{subfigure}{0.45\textwidth}
        \centering
        \includegraphics[width=\linewidth]{2-срез.png}
        \caption{Первое сечение поперек дорожек}
        \label{fig:cross_section1}
    \end{subfigure}
    \hfill
    \begin{subfigure}{0.45\textwidth}
        \centering
        \includegraphics[width=\linewidth]{image4.png}
        \caption{Второе сечение поперек дорожек}
        \label{fig:cross_section2}
    \end{subfigure}
    \caption{Сечения CD-диска на 1-й установке}
    \label{fig:cd_cross_sections}
\end{figure}

Получаем глубину пита из графика, изображенного на рис \ref{fig:cross_section2}, $h = 35$ нм, расстояние между дорожками из рис. \ref{fig:cd_surface}, $d = 0.6$ мкм, длину и ширину пит нельзя померить из-за неразличимости структуры.

На второй установке был произведен тот же эксперимент, рассмотрим результаты измерений.

\begin{figure}[H]
    \centering
    \includegraphics[width=0.8\linewidth]{relief.png}
    \caption{Рельеф CD-диска, полученный на 2-й установке}
    \label{fig:cd_relief}
\end{figure}

\begin{figure}[H]
    \centering
    \begin{subfigure}{0.45\textwidth}
        \centering
        \includegraphics[width=\linewidth]{sech1.png}
        \caption{Сечение вдоль дорожек}
        \label{fig:longitudinal_section}
    \end{subfigure}
    \hfill
    \begin{subfigure}{0.45\textwidth}
        \centering
        \includegraphics[width=\linewidth]{image7.png}
        \caption{Сечение поперек дорожек}
        \label{fig:transverse_section}
    \end{subfigure}
    \caption{Сечения CD-диска на 2-й установке}
    \label{fig:cd_sections_setup2}
\end{figure}

Длина пита на диске варьируется от 0,4 до 0,8 мкм, расстояние между питами - около 0,6 мкм (из сечения вдоль дорожек). Расстояние между дорожками $l = 0,5$ мкм, ширина пита $a = 0,3$ мкм. Глубина пита около 30 нм.

Для определения радиуса скругления зонда мы использовали следующую формулу: 
$$d = \sqrt{8R\Delta Z}$$ 
где $d = 0,7$ мкм мы уже нашли как период расположения пит, $\Delta Z = 30$ нм, получили $R \approx 20,4$ нм.

\subsection{Спектроскопия CD-диска}

Перейдем к спектроскопии CD-диска. Для первой установки спектры представлены на рисунках \ref{fig:spectra_set1}. Для второй установки спектры - на рисунках \ref{fig:spectra_set2}.

\begin{figure}[H]
    \centering
    \begin{subfigure}{0.3\textwidth}
        \centering
        \includegraphics[width=\linewidth]{спектр-1.png}
        \caption{Точка 1}
    \end{subfigure}
    \hfill
    \begin{subfigure}{0.3\textwidth}
        \centering
        \includegraphics[width=\linewidth]{спектр-2.png}
        \caption{Точка 2}
    \end{subfigure}
    \hfill
    \begin{subfigure}{0.3\textwidth}
        \centering
        \includegraphics[width=\linewidth]{спектр-3.png}
        \caption{Точка 3}
    \end{subfigure}
    \\
    \begin{subfigure}{0.3\textwidth}
        \centering
        \includegraphics[width=\linewidth]{спектр-4.png}
        \caption{Точка 4}
    \end{subfigure}
    \hfill
    \begin{subfigure}{0.3\textwidth}
        \centering
        \includegraphics[width=\linewidth]{спектр-5.png}
        \caption{Точка 5}
    \end{subfigure}
    \hfill
    \begin{subfigure}{0.3\textwidth}
        \centering
        \includegraphics[width=\linewidth]{спектр-6.png}
        \caption{Точка 6}
    \end{subfigure}
    \caption{Результаты спектроскопии на 1-й установке}
    \label{fig:spectra_set1}
\end{figure}

\begin{figure}[H]
    \centering
    \begin{subfigure}{0.45\textwidth}
        \centering
        \includegraphics[width=\linewidth]{спееектр-1.png}
        \caption{Точка 1 (в яме)}
    \end{subfigure}
    \hfill
    \begin{subfigure}{0.45\textwidth}
        \centering
        \includegraphics[width=\linewidth]{666.png}
        \caption{Точка 2 (на поверхности)}
    \end{subfigure}
    \caption{Результаты спектроскопии на 2-й установке}
    \label{fig:spectra_set2}
\end{figure}

\begin{figure}[H]
    \centering
    \includegraphics[width=0.8\linewidth]{картаа.png}
    \caption{Карта расположения точек, в которых проведена спектроскопия}
    \label{fig:spectroscopy_map}
\end{figure}

Рассмотрение обеих серий эксперимента позволяет заметить два характерных эффекта. Во-первых, вблизи точки $z = 0$ наклон кривой, снятой внутри пита, оказывается более пологим по сравнению с кривой, полученной на ровной поверхности. Во-вторых, обнаруживается гистерезис - расхождение между кривыми подвода (красная) и отвода (синяя) зонда.

Физическая природа первого эффекта связана с изменением градиента сил взаимодействия. Более крутой наклон соответствует резкому изменению ван-дер-ваальсовых сил при приближении зонда. Находясь внутри пита, зонд взаимодействует не только с дном, но и с боковыми стенками, что уменьшает эффективное расстояние до материала и приводит к более раннему возникновению нелинейных эффектов, проявляющемуся в растягивании зависимости амплитуды от расстояния.

Гистерезисное расхождение кривых объясняется необходимостью преодоления дополнительных сил адгезии при отводе зонда. В атмосферных условиях основную роль играют капиллярные силы, связанные с наличием адсорбционного слоя жидкости. Эти силы удерживают зонд в контакте с поверхностью, вызывая запаздывание восстановления амплитуды колебаний.

Аналогия может быть проведена с различием между силой трения покоя и трения скольжения: сила, требуемая для отрыва зонда от поверхности, превышает силу, необходимую для поддержания его в отведенном состоянии. Величина гистерезиса составляет примерно 10\%, что приводит к сдвигу областей сканирования и требует проведения измерений только при одном направлении хода сканера для исключения искажений.

\subsection{Исследование калибровочной решетки TGT01}

\subsubsection{Установка №1}

На этой установке для работы с тестовой решеткой TGT01 был использован зонд с другими параметрами.

\begin{figure}[H]
    \centering
    \begin{subfigure}{0.45\textwidth}
        \centering
        \includegraphics[width=\linewidth]{ачх_авто_2.png}
        \caption{Режим "авто"}
    \end{subfigure}
    \hfill
    \begin{subfigure}{0.45\textwidth}
        \centering
        \includegraphics[width=\linewidth]{ачх_ручной_2.png}
        \caption{Режим "ручной"}
    \end{subfigure}
    \caption{АЧХ зонда на установке №1}
    \label{fig:afc_setup1}
\end{figure}

Отсюда резонансная частота зонда $\nu = 7,24$ кГц и добротность $Q = 33,4$. Такая добротность является хорошим показателем качества зонда.

Период тестовой решетки известен и равен 2,12 мкм. В работе он измерялся с помощью АСМ следующим образом: была выставлена скорость сканирования 6000 нм/c и наблюдалась область размером 5×2,5 мкм.

\begin{figure}[H]
    \centering
    \includegraphics[width=1\linewidth]{сеченияяя.png}
    \caption{Вид сверху, сечения и их графики (установка №1)}
    \label{fig:grating_setup1}
\end{figure}

По периодической структуре вида сверху можно сделать вывод, что период решетки примерно равен 2 мкм. Проанализировав графики сечений, получим приблизительно $1,9 \pm 0,1$ мкм для каждого из них. Более точно измерить достаточно сложно в силу не очень четких данных на установке №1 и неудачного выбора мест для сечений. 

\begin{table}[H]
\centering
\begin{tabular}{|c|c|}
\hline
\textbf{Способ} & \textbf{h, нм} \\ 
\hline
Сечение 1 & 350 \\ 
\hline
Сечение 2 & 325 \\ 
\hline
Среднее & $337,5 \pm 12,5$ \\ 
\hline
\end{tabular}
\caption{Измерение высоты иглы h по 2 сечениям (установка №1)}
\label{tab:height_setup1}
\end{table}

\subsubsection{Установка №2}

На установке №2 аналогичная методика измерений, однако используется зонд, отличный от используемого в установке №1.

\begin{figure}[H]
    \centering
    \includegraphics[width=0.8\linewidth]{ачх_2.png}
    \caption{АЧХ в режиме "ручной" (установка №2)}
    \label{fig:afc_setup2}
\end{figure}

Отсюда резонансная частота $\nu = 8,73$ кГц и добротность $Q = 26,7$. Значение добротности ниже, чем в предыдущем опыте, но все еще высокое.

Скорость сканирования в данном случае 5000 нм/c, размер области наблюдения -- 7×7 мкм.


\begin{figure}[H]
    \centering
    \includegraphics[width=0.8\linewidth]{сеч_1_2.png}
    \caption{Сечение 1 (установка №2)}
    \label{fig:grating_sections_setup1}
\end{figure}

\begin{figure}[H]
    \centering
    \includegraphics[width=0.8\linewidth]{сеч_2_2.png}
    \caption{Сечение 2 (установка №2)}
    \label{fig:grating_sections_setup2}
\end{figure}



На виде сверху гораздо лучше видна периодическая структура решетки. По виду сверху можно определить, что период (по диагонали) примерно равен 1,5 мкм.

\begin{table}[H]
\centering
\begin{tabular}{|c|c|}
\hline
\textbf{Т (сеч.1), мкм} & \textbf{Т (сеч.2), мкм} \\
\hline
1,38 & 1,38 \\
\hline
1,25 & 1,42 \\
\hline
1,50 & 1,30 \\
\hline
1,25 & 1,38 \\
\hline
1,50 & 1,44 \\
\hline
1,50 &  \\
\hline
1,50 &  \\
\hline
1,38 &  \\
\hline
\end{tabular}
\caption{Измерение периода решетки (по диагонали) по 2 сечениям}
\label{tab:period_setup2}
\end{table}

Итого для сечений 1 и 2 получены соответственно $T_1 = 1,41 \pm 0,32$ мкм, $T_2 = 1,38 \pm 0,23$ мкм.

\begin{table}[H]
\centering
\begin{tabular}{|c|c|}
\hline
\textbf{h(сеч.1), нм} & \textbf{h(сеч.2), нм} \\
\hline
388 & 288 \\
\hline
325 & 375 \\
\hline
338 & 375 \\
\hline
350 & 250 \\
\hline
325 &  \\
\hline
\end{tabular}
\caption{Измерение высоты иглы h по 2 сечениям (установка №2)}
\label{tab:height_setup2}
\end{table}

Итого получаем $h_1 = 345,1 \pm 13,5$ нм, $h_2 = 322,0 \pm 14,8$ нм. 

Таким образом, установка №2 позволила сделать более точные измерения, основанные на большем количестве экспериментальных данных, однако значения для обеих установок достаточно близкие и показывают схожие тенденции.

\section{Вывод}

В результате выполнения лабораторной работы были успешно достигнуты все поставленные цели:

\begin{enumerate}
\item \textbf{Ознакомление с физическими принципами работы АСМ:} изучены основные принципы атомно-силовой микроскопии, включая типы силовых взаимодействий (короткодействующие и дальнодействующие силы), режимы работы (контактный, полуконтактный, бесконтактный) и устройство основных компонентов микроскопа.

\item \textbf{Исследование топографии игольчатой подложки:} 
\begin{itemize}
\item Для тестовой решетки TGT01 измерен период структуры: на установке №1 получено $1,9 \pm 0,1$ мкм, на установке №2 -- $1,41 \pm 0,32$ мкм и $1,38 \pm 0,23$ мкм
\item Измерена высота игл: на установке №1 -- $337,5 \pm 12,5$ нм, на установке №2 -- $345,1 \pm 13,5$ нм и $322,0 \pm 14,8$ нм
\item Полученные значения соответствуют ожидаемым параметрам решетки
\end{itemize}

\item \textbf{Исследование CD-диска:}
\begin{itemize}
\item Определены геометрические параметры пит: глубина 30-35 нм, ширина 0,3-0,6 мкм, расстояние между дорожками 0,5-0,6 мкм
\item Проведена спектроскопия поверхности, выявлены эффекты гистерезиса (∼10\%) и изменения градиента сил взаимодействия
\item Обнаружено различие в поведении зонда внутри пита и на ровной поверхности
\end{itemize}

\item \textbf{Оценка радиуса кривизны зонда:} по результатам измерений CD-диска определен радиус скругления зонда $R \approx 20,4$ нм.
\end{enumerate}

Работа продемонстрировала возможности атомно-силовой микроскопии для исследования наноструктур с высоким разрешением и позволила получить практические навыки работы с современным исследовательским оборудованием.

\end{document}
